\chapter*{Kurzfassung} %*-Variante sorgt dafür, das Abstract nicht im Inhaltsverzeichnis auftaucht
\phantomsection\pdfbookmark{Kurzfassung}{kurzfassung}
Die Verpackungsmaschinen bei der Firma MULTIVAC bestehen aus mehreren, miteinander vernetzten Maschinen, die kontinuierlich Zustände, Ereignisse und Prozessdaten austauschen müssen, damit ein reibungsloses Zusammenspiel zwischen ihnen funktioniert. 
Dieser Austausch erfolgt über das Nachrichtenprotokoll MQTT (\acrlong{acr:MQTT}). Die Analyse der dabei entstehenden Nachrichtenströme gestaltet sich insbesondere bei komplexen Anlagen und hohen Nachrichtenvolumina schwierig. Bereits bestehende Werkzeuge ermöglichen zwar die Anzeige von Nachrichten, sind jedoch für eine systematische Fehlersuche und Diagnose nur eingeschränkt geeignet und setzen ein umfangreiches Expertenwissen für die Kommunikation zwischen den Maschinen voraus.

Ziel der vorliegenden Arbeit ist die  Entwicklung einer Überwachungs- und Diagnose-Applikation zur Analyse der MQTT-basierten Maschinenkommunikation in Verpackungslinien. Die Anwendung soll Nachrichtenströme transparent darstellen, dauerhaft speichern und eine effiziente Analyse großer Datenmengen ermöglichen. Zusätzlich sollen Kommunikations- und Prozessfehler automatisch erkannt sowie Funktionen zur Teilnehmerüberwachung und Nachrichtensimulation bereitgestellt werden.

Zunächst werden bestehende Monitoring-Werkzeuge hinsichtlich ihres Funktionsumfangs sowie ihrer Eignung für den geplanten Anwendungsfall untersucht. Anschließend wird eine prototypische Anwendung in Python entwickelt, um technische Anforderungen zu identifizieren und die Machbarkeit der Umsetzung zu untersuchen. Darauf aufbauend wird eine Softwarearchitektur auf Basis des MVVM-Entwurfsmusters konzipiert. Die finale Umsetzung erfolgt in der Programmiersprache C\# unter Verwendung von WPF für die Benutzeroberfläche, MQTTnet für die MQTT-Kommunikation sowie SQLite zur persistenten Datenspeicherung. Die Verarbeitung der Nachrichten erfolgt ereignisgesteuert und asynchron, um auch hohe Nachrichtenraten effizient verarbeiten zu können.

Im Rahmen dieser Arbeit wird eine funktionsfähige Monitoring- und Diagnose-Applikation entwickelt. Die Anwendung ermöglicht die dauerhafte Speicherung und Volltextsuche von MQTT-Nachrichten, die Überwachung aller Kommunikationsteilnehmer sowie die Analyse von Prozessabläufen wie Linienstart, Linienstopp und Artikelwechsel. Zusätzlich werden Mechanismen zur automatischen Erkennung von Struktur- und Prozessfehlern implementiert. Darüber hinaus werden Funktionen zum manuellen Versand von MQTT-Nachrichten bereitgestellt. Mithilfe der SQLite-Datenbank und eines Paging-Konzepts können auch große Mengen an Kommunikationsdaten effizient verwaltet und analysiert werden. Die Architektur ist modular aufgebaut und kann zukünftig um weitere Analyse- und Diagnosefunktionen erweitert werden.
\clearpage

\chapter*{Abstract} %*-Variante sorgt dafür, das Abstract nicht im Inhaltsverzeichnis auftaucht
\phantomsection\pdfbookmark{Abstract}{abstract}
Packaging lines produced at the company MULTIVAC consist of multiple interconnected machines that continuously exchange information regarding machine states, events and process data to ensure reliable and coordinated operation. This communication is realized using the Message Queuing Telemetry Transport (MQTT) protocol. The analysis of the resulting message streams becomes increasingly challenging in complex production lines and under high message volumes. Although existing monitoring tools provide the capability to display sent MQTT messages, systematic fault diagnosis and communication analyses are only possible to a limited extent and often require expert knowledge.

The objective of this thesis is to develop a working monitoring and diagnostic application for MQTT-based communication. The application is intended to provide a transparent visualization of message streams, support the persistent storage of communication data, and enable the efficient analysis of large message volumes. Additionally, the application should automatically detect errors in communication and processes and provide functionality for participant monitoring and message simulation.

Initially, existing tools for monitoring and displaying MQTT messages are evaluated to determine whether they fulfill the requirements of the intended use case. Subsequently, a prototype is developed in Python to assess the technical feasibility of the application and to identify relevant requirements and challenges. Based on the insights gained, a software architecture following the Model-View-ViewModel (MVVM) pattern is designed. The final application is implemented in C\# using Windows Presentation Foundation (WPF) for the graphical user interface, MQTTnet for MQTT communication, and SQLite for persistent data storage. Message processing is performed in an event-driven and asynchronous manner, enabling the efficient handling of high message rates and large volumes of communication data.

As a result of this project, a fully functional monitoring and diagnostic application is developed. The application enables the persistent storage and retrieval of MQTT messages through integrated search and filtering mechanisms. Furthermore, it provides functionality for monitoring communication participants and analysing production-line processes, including line start, line stop, article change and continuous article change. In addition, mechanisms for the automatic detection of structural and process-related errors are implemented, alongside functionality for the manual simulation and transmission of MQTT messages. By utilizing an SQLite database in combination with a paging concept, the application is capable of efficiently managing and analysing large volumes of communication data. The software architecture is designed in a modular manner, allowing future monitoring and analysis features to be integrated with minimal effort. This ensures both maintainability and extensibility of the application for future requirements.

\clearpage